% !Mode:: "TeX:UTF-8"

\begin{cabstract}
本文研究了跨尺度噪声与高效光纤接口耦合协同优化的中性原子量子接口问题。量子网络与量子计算的发展对高效光子-原子接口提出了迫切需求，而中性原子系统因其成熟技术和良好量子存储特性成为热门研究平台。然而，当前实验中存在光纤耦合效率偏低以及多尺度噪声（从微观的自发辐射到宏观的机械振动）对保真度影响等关键问题。本文建立了基于Lindblad主方程的跨尺度噪声理论模型，通过数值仿真对模型进行验证，并与实验数据进行详细对比，提出了针对光纤耦合及噪声管理的实验参数优化策略，为实现高保真量子接口提供了理论与实验依据。
\end{cabstract}

\begin{eabstract}
This thesis investigates the collaborative optimization of cross-scale noise and efficient fiber interface coupling in neutral atom quantum interfaces. The development of quantum networks and quantum computing has created an urgent demand for efficient photon-atom interfaces, with neutral atom systems becoming popular research platforms due to their mature technology and excellent quantum storage characteristics. However, key issues exist in current experiments, including low fiber coupling efficiency and the impact of multi-scale noise (from microscopic spontaneous emission to macroscopic mechanical vibrations) on fidelity. This thesis establishes a cross-scale noise theoretical model based on the Lindblad master equation, validates the model through numerical simulation, compares it with experimental data in detail, and proposes experimental parameter optimization strategies for fiber coupling and noise management, providing theoretical and experimental basis for achieving high-fidelity quantum interfaces.
\end{eabstract} 